
The power of these notions may be reflected by the following statement:
\begin{lemma}
\label{lem:superstandard characterize coefficients}
    Let $\lambda, \mu \vdash N$ be two partitions. Then:
    \begin{enumerate}
        \item $\Des(T_\lambda) = [N-1] \setminus \{\lambda_{1}', \lambda_{1}'+\lambda_{2}', \dots\}$.
        \item If $\Des(T) = [N-1] \setminus \{\lambda_{1}', \lambda_{1}'+\lambda_{2}', \dots\}$ for some $T \in \Syt(\mu)$, then $\mu' \ge \lambda'$. Furthermore, if $\mu = \lambda$ then $T = T_\lambda$.
    \end{enumerate}
\end{lemma}
    
    


\begin{proof}
    The first assertion is obvious, so let us focus on the second assertion.
    
    Let $T \in \Syt(\mu)$ be a tableau, and assume that $\Des(T) = \Des(T_\lambda)$ and $T \ne T_\lambda$. We aim to show that $\mu' > \lambda'$. Let us denote the first column that differs between $T$ and $T_\lambda$ by $i$. The $i$-th column of $T_\lambda$ contains the entries $s+1,\dots,s+\lambda_{i}'$, where $s = \lambda_1' + \dots + \lambda_{i-1}'$. To prove $\mu' > \lambda'$, it suffices to show that all these entries also appear in the $i$-th column of $T$.
    
    Assume, by contradiction, that some of these entries appear in other columns of $T$. Let $x$ be the minimal such entry, and let $j \ne i$ be the column of $T$ that contains $x$. Since $s+1$ appears in the $i$-th column of $T$, we may assume that $x > s+1$. If $j < i$, this contradicts the assumption that $T$ and $T_\lambda$ agree in the first $i-1$ columns. On the other hand, if $j > i$, then $x$ appears in the first row of $T$, since it is the minimal entry of the $j$-th column of $T$. Consequently, we obtain that $x-1 \notin \Des(T)$, while $x-1 \in \Des(T_\lambda)$, which contradicts the assumption that $\Des(T) = \Des(T_\lambda)$.
\end{proof}

Lemma~\ref{lem:superstandard characterize coefficients} associates the Young diagram of shape $\lambda$ with the set $\Des(T_\lambda)$. As we will see later, this association is powerful in analysing the Schur coefficients of symmetric sets.


Standard Young tableaux with at most $2$ rows have a slightly stronger property:
\begin{lemma}
\label{lem:superstandard two rows characterize coefficients}
    Let $\lambda = (N-k_1,k_1)$ and $\mu = (N-k_2,k_2)$ be two partitions of $N$ with at most two parts each. Then:
    \begin{enumerate}
        \item $\Des(T_\lambda) = \odds(k_1)$.
        \item If $\Des(T) = \odds(k_1)$ for some $T \in \Syt(\mu)$, then $k_1 = k_2$ and $T = T_\lambda$.
    \end{enumerate}
\end{lemma}

    

\begin{proof}
    The first assertion is obvious, so let us focus on the second assertion.

    Let $T \in \Syt(\mu)$ be a tableau, and suppose that $\Des(T) = \Des(T_\lambda)$. Given that both $T$ and $T_\lambda$ have two rows each, it suffices to show that $\row_2(T) = \row_2(T_\lambda)$. Since $\Des(T) = \{1,3,\dots,2k_1-1\}$, it follows that $\{1,3,\dots,2k_1-1\} \subseteq \row_1(T)$ and $\{2,4,\dots,2k_1\} \subseteq \row_2(T)$. Consequently, we have $2k_1+1 \in \row_1(T)$. As $T$ has no descents beyond index $2k_1-1$, all entries greater than $2k_1$ must appear in the first row. Therefore, we conclude that $\row_2(T) = \{2,4,\dots,2k_1\}$.
\end{proof}

Now we are ready to prove that if a set is symmetric with respect to a sparse statistic then it is Schur-positive:
\begin{lemma}
\label{lem:if two rows and symmetric then schur positive}
    Let $\A$ be a symmetric set with respect to a sparse statistic $D:\A \to 2^{[N-1]}$, and denote $n = \lfloor\frac{N}{2} \rfloor$. Then $\A$ is Schur-positive, and its Schur expansion is
    \[
        \Q_{D} (\A) = \sum_{k=0}^n |\A(D = \odds(k))|\; s_{N-k,k}.
    \]
\end{lemma}
\begin{proof}
    The set $\A$ is assumed to be symmetric, so by Theorem~\ref{thm:criterion schur positive young}, we have
    \begin{equation}\label{eqn:proof symmetric set with no consecutive indices is schur positive}
        \sum_{a \in \A} \boldsymbol{t}^{D(a)} = \sum_{\lambda \vdash N} c_\lambda \sum_{T \in \Syt(\lambda)} \boldsymbol{t}^{\Des(T)},
    \end{equation}
    where $c_\lambda$ are the Schur coefficients of $\A$. It suffices to show that if there exists $k$ such that $\lambda=(N-k,k)$ then $c_\lambda = |\A(D = \odds(k))|$, and otherwise $c_\lambda=0$.

    If $\A = \emptyset$, then $\Q_{D}(\A) = 0$, and the statement holds. Therefore, we may assume that $\A \ne \emptyset$, and consequently, there exists $c_\lambda \ne 0$ for some partition $\lambda \vdash N$. Let $\mu \vdash N$ be a partition with $c_\mu \ne 0$, maximal in the conjugate order (\ie such that $c_\lambda = 0$ for every $\lambda \vdash N$ with $\lambda' > \mu'$).
    
    Equation~\eqref{eqn:proof symmetric set with no consecutive indices is schur positive} implies that the equation
    \begin{equation}\label{eqn:proof symmetric set with no consecutive indices is schur positive - specific set}
        |\A(D = J)| = \sum_{\lambda \vdash N} c_\lambda \left|\Syt(\lambda)(\Des = J)\right|
    \end{equation}
    holds for all $J \subseteq [N-1]$. Let us find the right-hand side of Equation~\eqref{eqn:proof symmetric set with no consecutive indices is schur positive - specific set} when substituting $J = \Des(T_\mu)$. Let $T \in \Syt(\lambda)$ be a tableau with $\Des(T) = \Des(T_\mu)$. By Lemma~\ref{lem:superstandard characterize coefficients}, we have $T = T_\mu$ or $\lambda' > \mu'$. However, if $\lambda' > \mu'$ then $c_\lambda = 0$. Therefore, if $T \in \Syt(\lambda)$ has $\Des(T) = \Des(T_\mu)$ and $c_\lambda \ne 0$, then $T = T_\mu$. Consequently, substituting $J = \Des(T_\mu)$ into Equation~\eqref{eqn:proof symmetric set with no consecutive indices is schur positive - specific set}, we find that $|\A(D = \Des(T_\mu))| = c_\mu \neq 0$.

    We may conclude that there exists an element $a \in \A$ with $D(a) = \Des(T_\mu)$. The set $D(a)$ is sparse, so $\{1,2\} \nsubseteq D(a)$. Lemma~\ref{lem:superstandard characterize coefficients} implies that $\{1,2\} \subseteq \Des(T_\mu)$ whenever $\mu_1' >2$, so we may deduce that $\mu_1' \le 2$. Therefore, $c_\lambda = 0$ for every partition $\lambda$ with more than $2$ parts.
    
    Thus, we can reformulate Equation~\eqref{eqn:proof symmetric set with no consecutive indices is schur positive - specific set} as
    \begin{equation}\label{eqn:proof symmetric set with no consecutive indices is schur positive:equation with only two rows}
        |\A(D = J)| = \sum_{k=0}^n c_k \left|\Syt(N-k,k)(\Des = J)\right|.
    \end{equation}
    
    By Lemma~\ref{lem:superstandard two rows characterize coefficients}, since only partitions into at most two parts are involved in Equation~\eqref{eqn:proof symmetric set with no consecutive indices is schur positive:equation with only two rows}, substituting $J = \Des(T_{\lambda})$ for $\lambda = (N-k,k)$ yields $|\A(D = \odds(k))| = c_k$, as required.
\end{proof}

\begin{proof}[Second proof of Lemma~\ref{lem:if sparse and good sizes then schur positive two rows}]
    The lemma follows directly from Lemma~\ref{lem:sufficent condition symmetry two rows} together with Lemma~\ref{lem:if two rows and symmetric then schur positive}.
\end{proof}

